% 固定两个输入，比较阈值假设空间与区间假设空间能够实现的标签安排。
\begin{tikzpicture}[figure picture,foundation picture,x=1mm,y=1mm,font=\figurefont\footnotesize]
  \fill[FigurePaper] (0,0) rectangle (108,72);
  \node[font=\figurefont\footnotesize,text=FigureMuted] at (29,66) {目标标签};
  \node[font=\figurefont\footnotesize,text=FigureMuted] at (68,66) {阈值假设};
  \node[font=\figurefont\footnotesize,text=FigureMuted] at (94,66) {区间假设};
  \draw[FigureGrid] (4,61) -- (104,61);
  \draw[FigureGrid] (54,11) -- (54,72);
  \draw[FigureGrid] (82,11) -- (82,72);
  \foreach \yy/\a/\b in {54/0/0,43/0/1,32/1/0,21/1/1}{
    \draw[FigureMuted,line width=.5pt] (12,\yy) -- (36,\yy);
    \ifnum\a=1
      \fill[FigureBlue] (18,\yy) circle (1.7);
    \else
      \draw[FigureRed,line width=1.1pt,fill=FigurePaper] (18,\yy) circle (1.7);
    \fi
    \ifnum\b=1
      \fill[FigureBlue] (30,\yy) circle (1.7);
    \else
      \draw[FigureRed,line width=1.1pt,fill=FigurePaper] (30,\yy) circle (1.7);
    \fi
    \node[text=FigureInk] at (45,\yy) {$(\a,\b)$};
    \node[text=FigureBlue] at (94,\yy) {可实现};
  }
  \node[text=FigureBlue] at (68,54) {可实现};
  \node[text=FigureBlue] at (68,43) {可实现};
  \node[text=FigureRed] at (68,32) {无法实现};
  \node[text=FigureBlue] at (68,21) {可实现};
  \draw[FigureGrid] (4,14) -- (104,14);
  \node[font=\figurefont\footnotesize,text=FigureRed] at (68,6) {未打散 $C$};
  \node[font=\figurefont\footnotesize,text=FigureBlue] at (94,6) {打散 $C$};
\end{tikzpicture}%
